% !TeX root = ../../main.tex
% sections/app/appendix_dispersion_dynamics.tex

\chapter{価格分散の動学方程式の導出}
\label{chap:appendix_dispersion_dynamics}

本付録では、本文中で用いられる価格分散項\(\Delta_t^H\)が従う動学方程式を、カルボ・プライシングの仮定から厳密に導出する。

\section*{定理6：価格分散の動学}
\label{sec:appendix_dispersion_dynamics_theorem6}
カルボ・プライシングの仮定の下で、経済全体の価格分散項\(\Delta_t^H\)は、以下の再帰的な方程式に従う。
\begin{equation}
\Delta_t^H=(1-\xi_H)N\left(\frac{\tilde{p}_t^H}{p_t^H}\right)^{-\theta^H}+\xi_H\left(\frac{p_t^H}{p_{t-1}^H}\right)^{\theta^H}\Delta_{t-1}^H
\label{eq:appendix_dispersion_dynamics_theorem6_main}
\end{equation}

\section*{証明}
\label{sec:appendix_dispersion_dynamics_proof}

この定理を証明するため、価格分散項の定義式から出発し、価格改定を行う家計と行わない家計のグループに分けて計算を進める。

\subsection*{1.カルボ・プライシングによる総和の分割}
\label{sec:appendix_dispersion_dynamics_sub1}
価格分散項\(\Delta_t^H\)の定義は以下の通りである。
\begin{equation}
\Delta_t^H\equiv\sum_{h\in H}\left(\frac{p_t^h}{p_t^H}\right)^{-\theta^H}
\label{eq:appendix_dispersion_dynamics_sub1_definition}
\end{equation}
カルボ・プライシングの仮定に従い、時点\(t\)における全ての家計の集合\(H\)を、価格を改定する\((1-\xi_H)\)の割合の家計（改定組）と、価格を据え置く\(\xi_H\)の割合の家計（非改定組）に分割する。
\begin{equation}
\Delta_t^H=\underbrace{\sum_{h\in\text{改定組}}\left(\frac{p_t^h}{p_t^H}\right)^{-\theta^H}}_{\text{第1項：価格改定グループ}}+\underbrace{\sum_{h\in\text{非改定組}}\left(\frac{p_t^h}{p_t^H}\right)^{-\theta^H}}_{\text{第2項：価格非改定グループ}}
\label{eq:appendix_dispersion_dynamics_sub1_decomposition}
\end{equation}

\subsection*{2.価格改定グループの計算}
\label{sec:appendix_dispersion_dynamics_sub2}
価格改定の機会を得た家計は、全て同じ新しい最適価格\(\tilde{p}_t^H\)設定する。このグループに属する家計の数は\((1-\xi_H)N\)であるため、第1項の合計は以下のようになる。
\begin{equation}
\sum_{h\in\text{改定組}}\left(\frac{p_t^h}{p_t^H}\right)^{-\theta^H}=(1-\xi_H)N\left(\frac{\tilde{p}_t^H}{p_t^H}\right)^{-\theta^H}
\label{eq:appendix_dispersion_dynamics_sub2_revising_group}
\end{equation}

\subsection*{3.価格非改定グループの計算}
\label{sec:appendix_dispersion_dynamics_sub3}
価格を改定しない家計は、前期の価格\(p_{t-1}^h\)をそのまま維持するため、\(p_t^h=p_{t-1}^h\)である。したがって、第2項は以下のように変形できる。
\begin{align}
\sum_{h\in\text{非改定組}}\left(\frac{p_t^h}{p_t^H}\right)^{-\theta^H}&=\sum_{h\in\text{非改定組}}\left(\frac{p_{t-1}^h}{p_t^H}\right)^{-\theta^H} \nonumber \\
&=\sum_{h\in\text{非改定組}}\left(\frac{p_{t-1}^h}{p_{t-1}^H}\frac{p_{t-1}^H}{p_t^H}\right)^{-\theta^H} \nonumber \\
&=\left(\frac{p_t^H}{p_{t-1}^H}\right)^{\theta^H}\sum_{h\in\text{非改定組}}\left(\frac{p_{t-1}^h}{p_{t-1}^H}\right)^{-\theta^H}
\label{eq:appendix_dispersion_dynamics_sub3_derivation}
\end{align}
ここで、カルボ・プライシングの重要な仮定、すなわち「今期、価格を改定しない家計」は、「前期に存在した全家計の中から、割合\(\xi_H\)で無作為に選び出された標本（サンプル）である」という点を用いる。大数の法則によれば、家計の数\(N\)が十分に大きい場合、この無作為抽出されたサンプルの合計値は、母集団（前期の全家計）の合計値のほぼ\(\xi_H\)倍となる。
\begin{equation}
\underbrace{\sum_{h\in\text{非改定組}}\left(\frac{p_{t-1}^h}{p_{t-1}^H}\right)^{-\theta^H}}_{\text{サンプル（非改定組）の合計}}\approx\xi_H\times\underbrace{\sum_{h\in\text{全家計}}\left(\frac{p_{t-1}^h}{p_{t-1}^H}\right)^{-\theta^H}}_{\text{母集団（前期の全家計）の合計}}
\label{eq:appendix_dispersion_dynamics_sub3_sampling_assumption}
\end{equation}
右辺の「母集団の合計」は、まさしく前期の価格分散\(\Delta_{t-1}^H\)の定義そのものである。したがって、以下の近似が成立する。
\begin{equation}
\sum_{h\in\text{非改定組}}\left(\frac{p_{t-1}^h}{p_{t-1}^H}\right)^{-\theta^H}\approx\xi_H\Delta_{t-1}^H
\label{eq:appendix_dispersion_dynamics_sub3_sampling_result}
\end{equation}
これを代入すると、第2項は\(\xi_H(p_t^H/p_{t-1}^H)^{\theta^H}\Delta_{t-1}^H\)と近似できる。

\subsection*{4.方程式の再結合と結論}
\label{sec:appendix_dispersion_dynamics_sub4}
ステップ2とステップ3の結果を結合することで、定理で示された\(\Delta_t^H\)の動学方程式が得られる。
\begin{equation}
\Delta_t^H=(1-\xi_H)N\left(\frac{\tilde{p}_t^H}{p_t^H}\right)^{-\theta^H}+\xi_H\left(\frac{p_t^H}{p_{t-1}^H}\right)^{\theta^H}\Delta_{t-1}^H
\label{eq:appendix_dispersion_dynamics_sub4_final_form}
\end{equation}
（証明終）